The goal for this file (StoneCauchy.tex) is to give the equivalent in SSD of the following statement in classical mathematics: 
\begin{theorem}\label{GoalThmUniformCompleteExtension}
  If $Q\subseteq X$ has the subset uniformity of $X$ and is dense, and $f: Q \to Y$ is uniformly continuos and $Y$ complete, there exists an extension $X \to Y$. 
\end{theorem}
I have a proof where we first approach each element of $X$ with a sequence from $q$, and work with sequences instead. I feel that should not be the final version of the proof though, which should not work with sequences (I don't like it as it uses choice, which leads to some arbitrary mere existence statement one has to work around, which is not very elegant), but with general Cauchy filters. 

\begin{lemma}
  Let $S:\Stone$ and $x:\mathbb N \to S$ be such that %(relatively uniformly continuous w.r.t. \Noo)
  for all $E\subseteq S \times S$ entourage there merely exists some $N:\mathbb N$ such that whenever $n,m\geq N$, we have $(x_n,x_m) \in E$. Then we can construct an extension $y:\Noo \to S$. 
\end{lemma}



\begin{proof}
  We will create a map $2^S \to B_\infty$. 
  Let $D\subseteq S$ decidable. Consider $$f_D = \{n : \mathbb N | x_n \in D\}$$ We claim that $f_D$ is a finite or cofinite subset of $\mathbb N$, hence describes an element of $B_\infty$. 
  Then there merely exists a map $y : \Noo \to S$ extending $x$. 
  Consider
  $$
  \# D := (D \times D) \bigcup (\neg D \times \neg D)\subseteq S \times S
  $$
  Note that as $D$ is decidable, for each $s:S$ we have $(s,s)\in \# D$. Also $\#D$ is decidable, hence open, hence an entourage of $S$. 
  Therefore, by assumption there merely exists some $N:\mathbb N$ such that $n,m\geq N$ implies $x_n,x_m \in \#D$. Note that if $x_N \in D$, it follows for all $n>N$ that $(x_N,x_n) \in D\times D$, hence $x_n\in D$, thus $f_D$ is cofinite. 
  Similarly, if $x_N \in \neg D$, it follows that $f_D$ is finite. 
  As $D$ is decidable, we conclude that $f_D$ is indeed finite or cofinite.

  Hence $f$ defines a function $2^S \to B_\infty$. 
  Therefore it dualizes to a function $y:\Noo \to S$. 
  To show that $y$ extends $x$, we need to show that for each $n:\mathbb N$, when we apply $f$ to a set containing $x_n$, we get a set containing $n$, which is clear from the definition. 
\end{proof}

\begin{lemma}
  For any entourage $E$ of $\Noo$, there exists some $N:\mathbb N$ such that $\mathbb N_{\geq N} \times \mathbb N_{\geq N}\subseteq E$. 

\end{lemma} 
\begin{proof}
  By the basis theorem for products and the observation that $(\infty,\infty) \in E$, there must be some open subsets $U,V\subseteq \Noo$ such that $(\infty, \infty) \in U \times V \subseteq E$. And as the decidables form a basis for a Stone space, there are some decidable $D,E\subseteq \Noo$ containing $\infty$ with $D\subseteq U, E \subseteq V$. And as $\infty \in D,E$, both are cofinite, hence there exists an $N:\mathbb N$ with $\mathbb N_{\geq N} \subseteq D,E$. The required property follows for this $N$. 
\end{proof}

\begin{corollary}\label{BabyExtension}
  Suppose $S:\Stone$ and $x:\N \to S$ such that for all $E$ entourage of $S$, there is an entourage $F$ of $\Noo$ such that $(n,m)\in F$ implies $(x_n,x_m)\in E$. 
  Then we can construct an extension $\Noo \to S$. 
\end{corollary}
\begin{definition}
  We denote the maps $\N \to S$ with the above property a \textbf{Cauchy sequence}. If $x$ is such a Cauchy sequence, and $y$ the above extension, we call $y$ the \textbf{extension} of $x$, and $y(\infty)$ the \textbf{limit} of $x$. 
\end{definition}
\begin{lemma}
  Let $x: \N \to S$ be a Cauchy sequence and $L:S$. TFAE: 
  \begin{enumerate}[(i)]
    \item\label{Llimit} $x$ has the limit $L$. 
    \item\label{Lcofinite} For each $D:2^S$, the set $\{n : \N | x_n \in D\}$ is cofinite iff $L\in D$. 
    \item\label{Lcofinal} For each $D:2^S$ with $x\in D$, there exists an $N:\mathbb N$ such that if $n>N$, then $x_n \in D$. 
  \end{enumerate}
\end{lemma}
\begin{proof}
  Clearly \ref{Lcofinal} and \ref{Lcofinite} are equivalent. 
  For the converse direction, note that for $y:\Noo \to S$ dual to $f:2^S\to N_{(co)fin}$, 
  we have that $y(\infty) = s$ iff for all decidables TODO
\end{proof}

Note that \Cref{BabyExtension} is an instance of \Cref{GoalThmUniformCompleteExtension}. We will use it to prove a more general case. 

\begin{definition}
  \item Let $f:A \to B$. An \textbf{entourage of $A$ relative to $f$} is an entourage of $A$ that is of the form $(f\times f)^{-1} (E)$ for some entourage $E$ of $B$. 

  \item A function $g:B \to C$ is called \textbf{relatively uniformly continuous} if for all entourages $E$ of $C$, $(g\times g)^{-1}(E)$ is an entourage of $B$ relative to $C$. 
  \item A subset $Q\subseteq S$ is called \textbf{dense} iff for all $U\subseteq S$ open such that there merely exists some $s:S$ with $U(s)$, there mereley exists some $q:Q$ with $U(q)$. 
\end{definition}

\begin{lemma}
  Let $Q\subseteq S$ be a dense subset of a Stone space.
  Then for each $s:S$ there merely exists a sequence $q:\mathbb N \to Q$ whose limit is $s$. 
\end{lemma}
\begin{proof}
  Let $e:\mathbb N \twoheadrightarrow 2^S$ enumerate the decidable subsets of $S$. Consider the sequence of descending opens given by
  $$
  O_{s,n} = \bigcap_{k<n \\ s \in b_O(e_k)} b_O(e_k).
  $$
  Note that each of these countably many opens is inhabited by $s$. By density and countable choice, we can get a sequence $q_{s,n}:Q$ such that $q_{s,n} \in O_{s,n}$ for all $n:\mathbb N$. 
  Furthermore, let $D\subseteq S$ decidable be such that $x\in D$. 
  We claim that TODO

\end{proof}


\begin{theorem}
  Any relatively uniformly continuous function from an overtly discrete dense subset of a Stone space can be extended to a function. 
\end{theorem}
\begin{proof}
  Let $Q\subseteq S$ be overtly discrete dense and let $f:Q \to T$ be relatively uniformly continuous. 
  Fix an enumeration $e: \mathbb N \to 2^S$ of the basis of $S$. Let $s:S$. We will show that there exists some $t_s:T$ such that 
  \begin{itemize}
    \item For any $q: \Noo \to Q$ such that $q_n\in Q$ for all $n:\mathbb N$ and $q_\infty = s$, 
      we have that $f(q_n) : \mathbb N \to T$ extends to a function $\Noo \to T$ sending $\infty$ to $t_s$. 
  \end{itemize}

  Let $s:S$. Consider $\mathcal N_s = \{D:2^S | s \in D \}$, which is a decidable subset of an overtly discrete space, hence overtly discrete. It hence has an enumeration $e:\mathbb N \to \mathcal N_s$. 
\end{proof}







